\chapter{Отождествление дефектной связности с семейным расслоением}

Предыдущая глава доказала, что полная $SO(3)$-связность устраняет
директорное ядро. Но это ещё могло означать добавление нового поля после
обнаружения проблемы. Настоящий гейт проверяет более строгую возможность:
является ли эта связность уже построенной семейной связностью на общем
трёхпоколенческом пакете.

\section{Одно расслоение вместо двух связностей}

Амальгамированный физический угол Тома V имеет вид
\begin{equation}
 \mathcal H_{45}=E_{\rm fam}\otimes H_{15},
 \qquad \rank_{\mathbb R}E_{\rm fam}=3.
 \label{eq:v6-family-bundle-h45}
\end{equation}
Поле порядка является не независимым представлением, а функториальной
конструкцией из того же семейного расслоения:
\begin{equation}
 Q\in\Gamma\!\left(\operatorname{Sym}^2_0E_{\rm fam}\right).
 \label{eq:v6-Q-induced-family-bundle}
\end{equation}
Поэтому одна связность $\nabla^B=d+B$ на $E_{\rm fam}$ автоматически
индуцирует
\begin{align}
 \nabla^B(\psi)&=d\psi+B\psi,
 &\psi&\in\Gamma(E_{\rm fam}\otimes H_{15}),
 \label{eq:v6-family-fermion-connection}\\
 \nabla^B Q&=dQ+[B,Q],
 &Q&\in\Gamma(\operatorname{Sym}^2_0E_{\rm fam}).
 \label{eq:v6-family-Q-connection}
\end{align}

Явная матричная проверка индуцированного пятимерного представления дала
\begin{equation}
 -\sum_{a=1}^3L_a^2=6I_5,
\end{equation}
то есть точно прежний семейный спин $j=2$. Максимальный остаток между
матричной формулой $[B,Q]$ и действием индуцированных генераторов равен
$3.7\cdot10^{-15}$.

\begin{proposition}[Единственность типа связности]
Полная связность предыдущей главы и семейная связность локального
трубчатого переноса имеют один и тот же геометрический носитель
$E_{\rm fam}$. На фермионах она действует в триплете, а на $Q$ --- в
индуцированном симметричном бесследовом квадрате. Вторая независимая
$SO(3)$-связность не требуется.
\end{proposition}

Это отождествление не использует запрещённую калибровку всей алгебры
$M_{300}$. Выделяется только уже зафиксированный семейный фактор
$E_{\rm fam}$.

\section{Нормировка общим следом}

Пусть $F_B\in\Omega^2(so(3))$. На физическом углу его действие равно
\begin{equation}
 F_B\otimes I_{15}.
\end{equation}
Тогда нормированный след даёт точное тождество
\begin{equation}
 \frac1{45}\Tr_{45}
 \left((F_B^*F_B)\otimes I_{15}\right)
 =\frac13\Tr_3(F_B^*F_B).
 \label{eq:v6-family-connection-trace-normalization}
\end{equation}
На вещественно удвоенном углу ранга девяносто получается то же значение.
Для общего проверочного элемента оба остатка равны соответственно нулю и
$5.6\cdot10^{-17}$.

Следовательно, после задания внешнего пространственного исчисления общий
след уже нормирует кривизну семейной связности. Независимый относительный
вес между её действием на поколениях и на поле $Q$ не нужен: второе
действие индуцировано первым.

Однако сохраняется прежняя граница. След нормирует заданное вложение, но
не создаёт сам пространственный дифференциал и не доказывает, что
глобальная семейная симметрия обязана стать локальной динамической
симметрией.

\section{Аномальный аудит физического пакета}

Все пятнадцать левых вейлевских каналов одного стандарт-модельного
поколения объединяются в вещественный семейный триплет $j=1$. Для этого
представления симметризованный кубический след равен нулю:
\begin{equation}
 \Tr\bigl(t_a\{t_b,t_c\}\bigr)=0.
\end{equation}
Поэтому локальная кубическая $SO(3)^3$-аномалия отсутствует.

Единственный потенциально нетривиальный смешанный коэффициент с двумя
семейными генераторами пропорционален сумме гиперзарядов всех левых
вейлевских каналов:
\begin{equation}
 6\frac16+3\left(-\frac23\right)
 +3\frac13+2\left(-\frac12\right)+1=0.
 \label{eq:v6-family-hypercharge-anomaly-sum}
\end{equation}
Остальные смешанные коэффициенты с одним семейным генератором исчезают
из-за его нулевого следа.

Глобальная аномалия Виттена также отсутствует. Для целочисленного
представления $j=1$ удвоенный индекс Дынкина равен четырём, поэтому
\begin{equation}
 15\cdot4=60\equiv0\pmod2.
\end{equation}
Более новая аномалия $SU(2)$ относится к определённым полуцелым
представлениям и к текущему целочисленному $SO(3)$-пакету не применяется;
см. DOI \path{10.1016/0370-2693(82)90728-6} и
\path{arXiv:1810.00844}.

\begin{theorem}[Аномальная допустимость семейной калибровки]
Существующий физический пакет из пятнадцати стандарт-модельных каналов в
семейном триплете не имеет локальной $SO(3)^3$-аномалии, смешанной
$SO(3)^2U(1)_Y$-аномалии и глобальной аномалии Виттена. Новые фермионные
состояния для аномального сокращения не требуются.
\end{theorem}

Это отличается от закрытой замены $3\to2\oplus1$: там появлялись
пятнадцать фундаментальных дублетов и нечётная глобальная аномалия. Здесь
исходный целочисленный триплет не меняется.

\section{Старый тетраэдрический сектор и последний генератор}

Одноосное поле $Q$ нарушает $SO(3)$ только до $O(2)$ и оставляет один
непрерывный генератор. Но ранний слой проекта уже содержит канонический
тензор
\begin{equation}
 \mathcal T_{ijk}=\sum_{a=1}^4(n_a)_i(n_a)_j(n_a)_k,
 \qquad
 \operatorname{Stab}_{SO(3)}(\mathcal T)=A_4,
 \label{eq:v6-family-tetrahedral-tensor}
\end{equation}
построенный из четырёх проекторных осей без нового непрерывного
направления.

Массовая матрица связности от одного $Q$ имеет спектр
\begin{equation}
 \{0,1.5077157793\ldots,1.5077157793\ldots\}
\end{equation}
в используемой нормировке генераторов. Для канонического
$\mathcal T$ тензорный вклад имеет спектр
\begin{equation}
 \left\{\frac{128}{9},\frac{128}{9},\frac{128}{9}\right\}.
\end{equation}
Совместный спектр равен
\begin{equation}
 \{14.2222222222\ldots,
 15.7299380015\ldots,
 15.7299380015\ldots\}.
 \label{eq:v6-family-combined-gauge-mass-map}
\end{equation}
Непрерывного ядра больше нет. Если ось $Q$ совпадает с одной из вершин
тетраэдра, остаточная группа равна
\begin{equation}
 A_4\cap O(2)=\operatorname{Stab}_{A_4}(n_a)\simeq\mathbb Z_3.
\end{equation}

Это согласуется с явными калибровочными моделями, где поле спина три
нарушает $SO(3)$ до $A_4$ и поддерживает конечные вихри с голономиями
порядков два и три; см. \path{arXiv:0910.4392} и
\path{arXiv:2607.12366}.

\section{Последняя незакрытая граница}

Теперь четыре прежние проблемы разделены точнее:
\begin{enumerate}
 \item тип полной связности: \textbf{выведен из одного семейного
 расслоения};
 \item её следовая нормировка: \textbf{получена};
 \item аномальная допустимость: \textbf{получена};
 \item отсутствие непрерывного стабилизатора при $Q+\mathcal T$:
 \textbf{получено кинематически}.
\end{enumerate}

Но пока не доказано, что тетраэдрический тензор конденсируется из того же
родительского действия с ненулевой жёсткостью. Его групповой тип и
каноническая форма известны, однако коэффициент его кинетического и
потенциального вклада не выведен из $\tau_{300}$. Без этого нельзя
использовать положительный спектр~\eqref{eq:v6-family-combined-gauge-mass-map}
как физическую массу калибровочных бозонов.

Кроме того, внешний пространственный дифференциал и переход от глобального
семейного переноса к динамической локальной симметрии остаются отдельным
родительским условием. $M_{300}$ служит контейнером и мерой, но не
координатной калибровочной алгеброй.

\begin{nogocriterion}[Динамическая локализация ещё не выведена]
Нельзя объявлять семейные калибровочные бозоны физически полученными, пока
единое действие не порождает одновременно кривизну $F_B$, ковариантные
производные $Q$ и $\mathcal T$, ненулевой тетраэдрический конденсат и их
относительные коэффициенты без нового секторного веса.
\end{nogocriterion}

\section{Вердикт}

\begin{protocol}
\begin{enumerate}
 \item одна связность на фермионах и $Q$: \textbf{получена};
 \item индуцированное представление $Q$: \textbf{$j=2$};
 \item вторая независимая $SO(3)$-связность: \textbf{не требуется};
 \item нормировка на $H_{45}$ и $H_{90}$: \textbf{совпадает точно};
 \item локальные аномалии: \textbf{отсутствуют};
 \item глобальная аномалия Виттена: \textbf{отсутствует};
 \item новые фермионы для сокращения аномалий: \textbf{не требуются};
 \item непрерывный стабилизатор одного $Q$: \textbf{одномерный};
 \item непрерывный стабилизатор пары $Q+\mathcal T$: \textbf{нулевой};
 \item остаточная дискретная группа: \textbf{$\mathbb Z_3$};
 \item тетраэдрический вклад из того же действия: \textbf{не выведен};
 \item динамическая локализация семейной симметрии: \textbf{не выведена};
 \item рождение материи: \textbf{не закрыто};
 \item следующий гейт: \textbf{родительский вывод тетраэдрической массы
 полной семейной связности}.
\end{enumerate}
\end{protocol}

Машинный сертификат сохранён в
\path{s2t_v6_bosonic_defect_family_connection_parent_identification_gate_results.json}.